\documentclass[bookchapter]{unnes-ta_plus}
\usepackage[indonesian]{babel}
\usepackage{pdfpages}

\unnesSetTitle{JUDUL NASKAH CHAPTER ANDA}
\unnesSetAuthor{Nama Mahasiswa}
\unnesSetNIM{1234567890}
\unnesSetProdi{Nama Prodi}
\unnesSetFakultas{Nama Fakultas}
\unnesSetDegree{Sarjana ....}
\unnesSetCity{Semarang}
\unnesSetYear{2026}
\unnesSetPembimbing{Nama Pembimbing}{NIP Pembimbing}

% Khusus bookchapter:
\booktitle{Judul Book Chapter / Buku}
\bookpublisher{Nama Penerbit}
\bookstatus{akan/telah} % sesuaikan

\begin{document}

\maketitleunnes
\unnesPersetujuanPembimbing
% \unnesPengesahanPenguji{Ketua}{Sekretaris}{Penguji 1}{Penguji 2}
% \unnesPernyataanKeaslian

\tableofcontents
\unnesMainMatter

\chapter{Naskah Chapter}
% \includepdf[pages=-]{naskah_chapter.pdf}
Tuliskan naskah chapter di sini atau sisipkan PDF seperti contoh di atas.

\end{document}
